\documentclass{mpit}
\usepackage[utf8]{inputenc}
\usepackage[T1]{fontenc}
\usepackage{geometry}
\usepackage{graphicx}
\usepackage{float}
\usepackage{hyperref}
\usepackage[english,indonesian]{babel}

\title{Optimasi Pemrosesan Template LaTeX Menggunakan Pendekatan Berbasis AI}
\author{Academic Support Agent}
\date{\today}

\begin{document}
\maketitle

\begin{abstract}

\end{abstract}

\chapter{PENDAHULUAN}


\chapter{TINJAUAN PUSTAKA}


% Original content (commented out):
% %% ESTE TEMPLATE FOI ADAPTADO A PARTIR DA CLASSE ppgccufscar
%% criada pelo Programa de Pós Graduação em Ciência da Computação da UFSCAR
%%CRÉDITOS:Profa. Sandra Fabbri.
%% A classe ppgccufscar foi feita por Daniel Beck,
%% Daniel Bruno e Prof. Marcio


%%
%% Opcoes que podem ser passadas 'a classe
%%
%% Todas as opcoes da classe abnt (AbnTeX) sao validas.
%% Outras opcoes sao: quali e tese (dissertacao e' padrao)
%%

\documentclass[quali]{mpit}
%\documentclass{mpit}

%% pacotes que deseja usar
%% pacotes incompativeis sao:
%%   qualquer pacote de citacoes, como natbib, apalike, cite, etc.
%% o pacote babel ja vem carregado com ingles e portugues

\usepackage[utf8]{inputenc}
\usepackage{graphicx}

% Change these to change out the sort outputs.

\titulo{Exemplo de título grande para ocupar mais de uma linha}
\autor{Fulano de Tal}
\orientador[Orientadora]{Profa. Dra. Fulana de Tal}
\coorientador{Prof. Dr. Beltrano da Silva}
\areaconcentracao{Engenharia de Software}
\data{Mês/ano}

% epigrafre, agradecimentos sao feitos 'na mao'.
% um dia eu faco alguns comandos para eles ;)

\begin{document} 

\capa
\folhaderosto

\dedicatoria{A meus pais.}
\begin{agradecimentos}

\end{agradecimentos}
\epigrafe{Se você não consegue explicar algo de forma simples, você não entendeu suficientemente bem.}{Albert Einstein}

\begin{resumo}
Lorem ipsum dolor sit amet, consectetur adipiscing elit. Proin euismod risus sed velit congue vitae posuere risus condimentum. Aliquam interdum condimentum risus, vel aliquet nibh pretium at. Integer egestas, dolor eget ultrices gravida, turpis nunc suscipit erat, ut aliquam velit nisl in turpis. Nam vel sodales justo. Nulla auctor, quam sed porta hendrerit, mi tortor rutrum ipsum, sed ullamcorper nisl erat nec ipsum. Class aptent taciti sociosqu ad litora torquent per conubia nostra, per inceptos himenaeos. Etiam eget arcu id nisi auctor iaculis. Curabitur dignissim porta lacus. Curabitur bibendum aliquet diam vel tempor. Maecenas vulputate aliquet purus eget interdum. Etiam pretium tortor in mi interdum volutpat. Aenean at convallis lectus. Nulla adipiscing massa eget diam consequat a consequat nisi pharetra. Suspendisse potenti. Sed at tellus metus, at tempor turpis. Curabitur vitae ligula libero, et aliquet diam. Proin bibendum hendrerit interdum. Curabitur luctus, augue sit amet pharetra consectetur, erat tellus sagittis ligula, eget rhoncus neque diam mattis urna. Donec consequat pretium euismod. Nam pellentesque facilisis quam et scelerisque. Aenean eu lacus ac sem vehicula laoreet a quis erat. Nulla facilisi. Sed a tortor eros, et viverra felis. Duis molestie elit in arcu condimentum quis blandit justo fermentum. Integer rhoncus accumsan hendrerit.
\palavraschave{Latex, Template, Exemplo}
\end{resumo}

\begin{abstract}
Lorem ipsum dolor sit amet, consectetur adipiscing elit. Proin euismod risus sed velit congue vitae posuere risus condimentum. Aliquam interdum condimentum risus, vel aliquet nibh pretium at. Integer egestas, dolor eget ultrices gravida, turpis nunc suscipit erat, ut aliquam velit nisl in turpis. Nam vel sodales justo. Nulla auctor, quam sed porta hendrerit, mi tortor rutrum ipsum, sed ullamcorper nisl erat nec ipsum. Class aptent taciti sociosqu ad litora torquent per conubia nostra, per inceptos himenaeos. Etiam eget arcu id nisi auctor iaculis. Curabitur dignissim porta lacus. Curabitur bibendum aliquet diam vel tempor. Maecenas vulputate aliquet purus eget interdum. Etiam pretium tortor in mi interdum volutpat. Aenean at convallis lectus. Nulla adipiscing massa eget diam consequat a consequat nisi pharetra. Suspendisse potenti. Sed at tellus metus, at tempor turpis. Curabitur vitae ligula libero, et aliquet diam. Proin bibendum hendrerit interdum. Curabitur luctus, augue sit amet pharetra consectetur, erat tellus sagittis ligula, eget rhoncus neque diam mattis urna. Donec consequat pretium euismod. Nam pellentesque facilisis quam et scelerisque. Aenean eu lacus ac sem vehicula laoreet a quis erat. Nulla facilisi. Sed a tortor eros, et viverra felis. Duis molestie elit in arcu condimentum quis blandit justo fermentum. Integer rhoncus accumsan hendrerit.
\keywords{Latex, Template, Example}
\end{abstract}

\listoffigures
\listoftables

%% sumario
\tableofcontents

%% aqui comeca o texto da monografia
%% voce pode dividir o conteudo em varios arquivos.
%% por exemplo, intro.tex, fundamentacao.tex, desenvolvimento.tex, conclusao.tex.
%% dai, vc inclui aqui assim: \input{intro.tex} e assim por diante.

\chapter{Introdução}

O capítulo de introdução deve contextualizar o trabalho, apresentar de forma clara os objetivos e finalizar com uma descrição de como o texto está estruturado. A contextualização, os objetivos e a estrutração do texto devem compor subseções, conforme apresentado neste modelo.

\section{Contextualização}

Uma boa contextualização deve apresentar justificativas para o trabalho. Estas, por sua vez devem estar baseadas numa revisão de literatura e devem estar devidamente referenciadas. A seguir, seguem alguns exemplos de citação.

Exemplo de citação indireta \cite{tolvanenKelly08}.
Exemplo de citação de autor: Segundo \citeonline{lucredio2009}, bla bla bla.
Exemplo de citação longa:

\begin{citacao}
A teleconferência permite ao indivíduo participar de um encontro
nacional ou regional sem a necessidade de deixar seu local de
origem. Tipos comuns de teleconferênia incluem o uso da televisão,
telefone e computador. Através de áudio conferência, utilizando
a companhia local de telefone, um sinal de áudio pode ser emitido
em um salão de qualquer dimensão \cite[p. 181]{weberZhang96}.
\end{citacao}

Algum texto...Lorem ipsum dolor sit amet, consectetur adipiscing elit. Cras at dui ligula. Vivamus eget est nulla. Suspendisse ligula elit, tempor ut facilisis vitae, mattis vitae ligula. Curabitur facilisis tempor nunc, eu vehicula neque tristique non. Aenean pellentesque odio non leo semper fringilla. Morbi semper, libero eu scelerisque egestas, tortor leo adipiscing lectus, quis molestie odio odio eu augue. Etiam mattis mi eu odio pellentesque imperdiet. Integer ac blandit elit. Maecenas elementum nunc id enim pellentesque consectetur. Phasellus porttitor iaculis enim sed consequat. Mauris scelerisque lacus nec nibh tempor in mattis urna tempus. Integer condimentum hendrerit ante nec commodo. Mauris suscipit, arcu et faucibus pretium, purus velit egestas purus, nec venenatis elit mauris nec ligula. Praesent lectus risus, consequat sit amet rutrum et, cursus in tellus.

\subsection{Subseção de exemplo}

Algum texto...Lorem ipsum dolor sit amet, consectetur adipiscing elit. Cras at dui ligula. Vivamus eget est nulla. Suspendisse ligula elit, tempor ut facilisis vitae, mattis vitae ligula. Curabitur facilisis tempor nunc, eu vehicula neque tristique non. Aenean pellentesque odio non leo semper fringilla. Morbi semper, libero eu scelerisque egestas, tortor leo adipiscing lectus, quis molestie odio odio eu augue. Etiam mattis mi eu odio pellentesque imperdiet. Integer ac blandit elit. Maecenas elementum nunc id enim pellentesque consectetur. Phasellus porttitor iaculis enim sed consequat. Mauris scelerisque lacus nec nibh tempor in mattis urna tempus. Integer condimentum hendrerit ante nec commodo. Mauris suscipit, arcu et faucibus pretium, purus velit egestas purus, nec venenatis elit mauris nec ligula. Praesent lectus risus, consequat sit amet rutrum et, cursus in tellus.

\subsubsection{Subsubseção de exemplo}

Algum texto...Lorem ipsum dolor sit amet, consectetur adipiscing elit. Cras at dui ligula. Vivamus eget est nulla. Suspendisse ligula elit, tempor ut facilisis vitae, mattis vitae ligula. Curabitur facilisis tempor nunc, eu vehicula neque tristique non. Aenean pellentesque odio non leo semper fringilla. Morbi semper, libero eu scelerisque egestas, tortor leo adipiscing lectus, quis molestie odio odio eu augue. Etiam mattis mi eu odio pellentesque imperdiet. Integer ac blandit elit. Maecenas elementum nunc id enim pellentesque consectetur. Phasellus porttitor iaculis enim sed consequat. Mauris scelerisque lacus nec nibh tempor in mattis urna tempus. Integer condimentum hendrerit ante nec commodo. Mauris suscipit, arcu et faucibus pretium, purus velit egestas purus, nec venenatis elit mauris nec ligula. Praesent lectus risus, consequat sit amet rutrum et, cursus in tellus.

\section{Objetivos}

Esta subseção deve apresentar os objetivos do trabalho: geral e específicos.

Algum texto...Lorem ipsum dolor sit amet, consectetur adipiscing elit. Sed consequat risus a enim iaculis et posuere metus adipiscing. Aliquam risus erat, elementum ut placerat vel, posuere vitae augue. Proin ac mattis ligula. Sed metus lectus, venenatis ac porta tristique, placerat non quam. Sed sit amet sem eu massa euismod condimentum. Praesent eget nulla id ligula luctus ultricies sit amet eu massa. Curabitur bibendum purus non ante tempus gravida eleifend urna aliquet. Fusce dolor nisl, ullamcorper nec dignissim id, ultrices id lectus. Proin auctor eros vel nunc lobortis ut pretium lectus tempus. Vestibulum lacinia gravida eros, quis tempus odio dapibus eu. Cras lobortis odio augue, eu molestie massa. Proin ac leo leo, ac lobortis est. Maecenas et aliquet sapien. Proin tempus viverra dui id lobortis.

\section{Organização do texto}

Esta subseção deve apresentar a estruturação do texto, descrevendo brevemente ao leitor o que poderá ser encontrado nos capítulos subsequentes e como eles se relacionam.

Algum texto... Lorem ipsum dolor sit amet, consectetur adipiscing elit. Proin euismod risus sed velit congue vitae posuere risus condimentum. Aliquam interdum condimentum risus, vel aliquet nibh pretium at. Integer egestas, dolor eget ultrices gravida, turpis nunc suscipit erat, ut aliquam velit nisl in turpis. Nam vel sodales justo. Nulla auctor, quam sed porta hendrerit, mi tortor rutrum ipsum, sed ullamcorper nisl erat nec ipsum. Class aptent taciti sociosqu ad litora torquent per conubia nostra, per inceptos himenaeos. Etiam eget arcu id nisi auctor iaculis. Curabitur dignissim porta lacus. Curabitur bibendum aliquet diam vel tempor. Maecenas vulputate aliquet purus eget interdum. Etiam pretium tortor in mi interdum volutpat.

Aenean at convallis lectus. Nulla adipiscing massa eget diam consequat a consequat nisi pharetra. Suspendisse potenti. Sed at tellus metus, at tempor turpis. Curabitur vitae ligula libero, et aliquet diam. Proin bibendum hendrerit interdum. Curabitur luctus, augue sit amet pharetra consectetur, erat tellus sagittis ligula, eget rhoncus neque diam mattis urna. Donec consequat pretium euismod. Nam pellentesque facilisis quam et scelerisque. Aenean eu lacus ac sem vehicula laoreet a quis erat.

\chapter{Revisão bibliográfica}

A revisão bibliográfica deve ser composta de duas subseções: fundamentação teórica e revisão sistemática.

Algum texto...Lorem ipsum dolor sit amet, consectetur adipiscing elit. Proin euismod risus sed velit congue vitae posuere risus condimentum. Aliquam interdum condimentum risus, vel aliquet nibh pretium at. Integer egestas, dolor eget ultrices gravida, turpis nunc suscipit erat, ut aliquam velit nisl in turpis. Nam vel sodales justo. Nulla auctor, quam sed porta hendrerit, mi tortor rutrum ipsum, sed ullamcorper nisl erat nec ipsum. Class aptent taciti sociosqu ad litora torquent per conubia nostra, per inceptos himenaeos. Etiam eget arcu id nisi auctor iaculis. Curabitur dignissim porta lacus. Curabitur bibendum aliquet diam vel tempor. Maecenas vulputate aliquet purus eget interdum. Etiam pretium tortor in mi interdum volutpat.

\section{Fundamentação teórica}

Um exemplo de figura é apresentado na figura \ref{calvin}.

%FIGURA: fluxo de processamento 
\begin{figure}[!htb]
\centering
\includegraphics [width=3.5in]{calvin}
\caption{Exemplo de figura}
\label{calvin}
\end{figure}

Algum texto...Lorem ipsum dolor sit amet, consectetur adipiscing elit. Proin euismod risus sed velit congue vitae posuere risus condimentum. Aliquam interdum condimentum risus, vel aliquet nibh pretium at. Integer egestas, dolor eget ultrices gravida, turpis nunc suscipit erat, ut aliquam velit nisl in turpis. Nam vel sodales justo. Nulla auctor, quam sed porta hendrerit, mi tortor rutrum ipsum, sed ullamcorper nisl erat nec ipsum. Class aptent taciti sociosqu ad litora torquent per conubia nostra, per inceptos himenaeos. Etiam eget arcu id nisi auctor iaculis. Curabitur dignissim porta lacus. Curabitur bibendum aliquet diam vel tempor. Maecenas vulputate aliquet purus eget interdum. Etiam pretium tortor in mi interdum volutpat.



Algum texto...Lorem ipsum dolor sit amet, consectetur adipiscing elit. Proin euismod risus sed velit congue vitae posuere risus condimentum. Aliquam interdum condimentum risus, vel aliquet nibh pretium at. Integer egestas, dolor eget ultrices gravida, turpis nunc suscipit erat, ut aliquam velit nisl in turpis. Nam vel sodales justo. Nulla auctor, quam sed porta hendrerit, mi tortor rutrum ipsum, sed ullamcorper nisl erat nec ipsum. Class aptent taciti sociosqu ad litora torquent per conubia nostra, per inceptos himenaeos. Etiam eget arcu id nisi auctor iaculis. Curabitur dignissim porta lacus. Curabitur bibendum aliquet diam vel tempor. Maecenas vulputate aliquet purus eget interdum. Etiam pretium tortor in mi interdum volutpat.

\section{Revisão sistemática}

Algum texto...Lorem ipsum dolor sit amet, consectetur adipiscing elit. Proin euismod risus sed velit congue vitae posuere risus condimentum. Aliquam interdum condimentum risus, vel aliquet nibh pretium at. Integer egestas, dolor eget ultrices gravida, turpis nunc suscipit erat, ut aliquam velit nisl in turpis. Nam vel sodales justo. Nulla auctor, quam sed porta hendrerit, mi tortor rutrum ipsum, sed ullamcorper nisl erat nec ipsum. Class aptent taciti sociosqu ad litora torquent per conubia nostra, per inceptos himenaeos. Etiam eget arcu id nisi auctor iaculis. Curabitur dignissim porta lacus. Curabitur bibendum aliquet diam vel tempor. Maecenas vulputate aliquet purus eget interdum. Etiam pretium tortor in mi interdum volutpat.

Algum texto...Lorem ipsum dolor sit amet, consectetur adipiscing elit. Proin euismod risus sed velit congue vitae posuere risus condimentum. Aliquam interdum condimentum risus, vel aliquet nibh pretium at. Integer egestas, dolor eget ultrices gravida, turpis nunc suscipit erat, ut aliquam velit nisl in turpis. Nam vel sodales justo. Nulla auctor, quam sed porta hendrerit, mi tortor rutrum ipsum, sed ullamcorper nisl erat nec ipsum. Class aptent taciti sociosqu ad litora torquent per conubia nostra, per inceptos himenaeos. Etiam eget arcu id nisi auctor iaculis. Curabitur dignissim porta lacus. Curabitur bibendum aliquet diam vel tempor. Maecenas vulputate aliquet purus eget interdum. Etiam pretium tortor in mi interdum volutpat.

\chapter{Proposta de Dissertação}

A proposta de dissertação deve estar estruturada da seguinte forma:
\begin{itemize}
\item questões de pesquisa;
\item metodologia;
\item resultados parciais e discussão
\item cronograma até a defesa (próximos passos)
\end{itemize}

Os seguintes subcapítulos ilustram esta organização.

\section{Questões de pesquisa}

Algum texto...Lorem ipsum dolor sit amet, consectetur adipiscing elit. Proin euismod risus sed velit congue vitae posuere risus condimentum. Aliquam interdum condimentum risus, vel aliquet nibh pretium at. Integer egestas, dolor eget ultrices gravida, turpis nunc suscipit erat, ut aliquam velit nisl in turpis. Nam vel sodales justo. Nulla auctor, quam sed porta hendrerit, mi tortor rutrum ipsum, sed ullamcorper nisl erat nec ipsum. Class aptent taciti sociosqu ad litora torquent per conubia nostra, per inceptos himenaeos. Etiam eget arcu id nisi auctor iaculis. Curabitur dignissim porta lacus. Curabitur bibendum aliquet diam vel tempor. Maecenas vulputate aliquet purus eget interdum. Etiam pretium tortor in mi interdum volutpat.

\section{Metodologia}

Algum texto...Lorem ipsum dolor sit amet, consectetur adipiscing elit. Proin euismod risus sed velit congue vitae posuere risus condimentum. Aliquam interdum condimentum risus, vel aliquet nibh pretium at. Integer egestas, dolor eget ultrices gravida, turpis nunc suscipit erat, ut aliquam velit nisl in turpis. Nam vel sodales justo. Nulla auctor, quam sed porta hendrerit, mi tortor rutrum ipsum, sed ullamcorper nisl erat nec ipsum. Class aptent taciti sociosqu ad litora torquent per conubia nostra, per inceptos himenaeos. Etiam eget arcu id nisi auctor iaculis. Curabitur dignissim porta lacus. Curabitur bibendum aliquet diam vel tempor. Maecenas vulputate aliquet purus eget interdum. Etiam pretium tortor in mi interdum volutpat.

\section{Resultados parciais e discussão}

Algum texto...Lorem ipsum dolor sit amet, consectetur adipiscing elit. Proin euismod risus sed velit congue vitae posuere risus condimentum. Aliquam interdum condimentum risus, vel aliquet nibh pretium at. Integer egestas, dolor eget ultrices gravida, turpis nunc suscipit erat, ut aliquam velit nisl in turpis. Nam vel sodales justo. Nulla auctor, quam sed porta hendrerit, mi tortor rutrum ipsum, sed ullamcorper nisl erat nec ipsum. Class aptent taciti sociosqu ad litora torquent per conubia nostra, per inceptos himenaeos. Etiam eget arcu id nisi auctor iaculis. Curabitur dignissim porta lacus. Curabitur bibendum aliquet diam vel tempor. Maecenas vulputate aliquet purus eget interdum. Etiam pretium tortor in mi interdum volutpat.

\section{Cronograma até a defesa}

Relacione as atividades aos meses que faltam para a defesa, conforme modelo mostrado na tabela \ref{crono}.

\begin{table}[h]
\centering
\caption{Cronograma (modelo de tabela)}
\label{crono}
\begin{tabular}{l | l | l | l | l}
\hline
\textbf{Atividade/Mês} & \textbf{1} & \textbf{2} & \textbf{3} & ... \\ \hline \hline
Atividade 1   & X  &   &   &     \\ \hline
Atividade 2   &   & X  &   &     \\ \hline
Atividade ... &   &   &  X &     \\ \hline
\end{tabular}
\end{table}


\begin{table}[h]
\centering
\caption{Cronograma para condução da pesquisa do aluno - Mestrado Profissional em Inovação Tecnológica (24 meses)}
\label{crono}
\begin{tabular}{l | l | l | l | l | l | l | l | l | l | l | l | l | l | l | l | l | l | l | l | l | l | l | l | l }
\hline
\textbf{Atividade/Mês} & \textbf{3} & \textbf{4} & \textbf{5} & \textbf{6}  & \textbf{7} & \textbf{8} & \textbf{9} & \textbf{10} & \textbf{11} & \textbf{12} & \textbf{1} & \textbf{2} & \textbf{3} & \textbf{4} & \textbf{5} & \textbf{6} & \textbf{7} & \textbf{8} & \textbf{9} & \textbf{10} & \textbf{11} & \textbf{12} & \textbf{1} & \textbf{2}\\ \hline \hline
Cumprir créditos de disciplinas  &   &   &X   &X   &X   &X   &   &X   &X   &X   &X   &X   &   &   &   &   &  &   &   &  &   \\ \hline
Revisão da proposta de projeto & X  &   &   &   &   &   &   &   &   &   &   &   &   &   &   &   &  &   &   &  &   \\ \hline
Preparar protocolo da revisão sistemática& X  &   &   &   &   &   &   &   &   &   &   &   &   &   &   &   &  &   &   &  &   \\ \hline

Seleção dos estudos primários & X  &   &   &   &   &   &   &   &   &   &   &   &   &   &   &   &  &   &   &  &   \\ \hline

Fazer extração de dados & X  &   &   &   &   &   &   &   &   &   &   &   &   &   &   &   &  &   &   &  &   \\ \hline

Fazer análise dos dados& X  &   &   &   &   &   &   &   &   &   &   &   &   &   &   &   &  &   &   &  &   \\ \hline

Formatar paper e submeter para publicação & X  &   &   &   &   &   &   &   &   &   &   &   &   &   &   &   &  &   &   &  &   \\ \hline

Fundamentação teórica clássica& X  &   &   &   &   &   &   &   &   &   &   &   &   &   &   &   &  &   &   &  &   \\ \hline

Preparação e submissão de projeto para o Comitê de Ética da UNIFESP & X  &   &   &   &   &   &   &   &   &   &   &   &   &   &   &   &  &   &   &  &   \\ \hline

Design do Estudo de Caso & X  &   &   &   &   &   &   &   &   &   &   &   &   &   &   &   &  &   &   &  &   \\ \hline

Protocolo de validação do Estudo de Caso & X  &   &   &   &   &   &   &   &   &   &   &   &   &   &   &   &  &   &   &  &   \\ \hline

\textbf{Qualificação} & X  &   &   &   &   &   &   &   &   &   &   &   &   &   &   &   &  &   &   &  &   \\ \hline

Desenvolvimento do Estudo de Caso & X  &   &   &   &   &   &   &   &   &   &   &   &   &   &   &   &  &   &   &  &   \\ \hline

Validação do Estudo de Caso & X  &   &   &   &   &   &   &   &   &   &   &   &   &   &   &   &  &   &   &  &   \\ \hline

Análise dos resultados& X  &   &   &   &   &   &   &   &   &   &   &   &   &   &   &   &  &   &   &  &   \\ \hline

Formatar paper e submeter para publicação& X  &   &   &   &   &   &   &   &   &   &   &   &   &   &   &   &  &   &   &  &   \\ \hline

Escrita e formatação da dissertação & X  &   &   &   &   &   &   &   &   &   &   &   &   &   &   &   &  &   &   &  &   \\ \hline

\textbf{Defesa} & X  &   &   &   &   &   &   &   &   &   &   &   &   &   &   &   &  &   &   &  &   \\ \hline
\end{tabular}
\end{table}


\chapter{Considerações finais}

Algum texto...Lorem ipsum dolor sit amet, consectetur adipiscing elit. Proin euismod risus sed velit congue vitae posuere risus condimentum. Aliquam interdum condimentum risus, vel aliquet nibh pretium at. Integer egestas, dolor eget ultrices gravida, turpis nunc suscipit erat, ut aliquam velit nisl in turpis. Nam vel sodales justo. Nulla auctor, quam sed porta hendrerit, mi tortor rutrum ipsum, sed ullamcorper nisl erat nec ipsum. Class aptent taciti sociosqu ad litora torquent per conubia nostra, per inceptos himenaeos. Etiam eget arcu id nisi auctor iaculis. Curabitur dignissim porta lacus. Curabitur bibendum aliquet diam vel tempor. Maecenas vulputate aliquet purus eget interdum. Etiam pretium tortor in mi interdum volutpat.

%%-------------------------------------------------
%% Referências
%% coloque aqui o seu arquivo .bib
%% IMPORTANTE: nao use bibliographystyle!
%% o estilo ja vem definido.
\bibliography{references}

%% de acordo com o template, o glossario vem
%% depois das referencias e deve estar em ordem
%% alfabetica.
%% depois de muito esforco consegui fazer com que
%% o glossario ficasse em ordem alfabetica automaticamente.
%% ainda nao sei a escalabilidade do algoritmo :(

%% DICA: voce pode ir definindo os acronimos ao longo do texto.
%% Por exemplo, no capitulo 1, vc ta escrevendo:
%% Segundo Fulano, Model-Driven Development (MDD)\acronym{MDD}{Model-Driven Development} � uma t�cnica bla bla bla...
\acronym{MDD}{Model-Driven Development}
\acronym{GPL}{General Purpose Language}
\acronym{SE}{Sistemas Embarcados}
\listofacronyms

\apendice

\chapter{Explicação do que é apêndice}

\begin{resumocap}
Dá pra usar o resuminho aqui também.
\end{resumocap}

Segundo a norma, apêndice é um documento elaborado pelo autor a fim de complementar sua argumentação sem afetar a unidade nuclear do trabalho.

\chapter{Outro apêndice}

Só pra ilustrar mais de um apêndice.

\anexo

\chapter{Explicação do que é anexo}

Anexos consistem em textos ou documentos não elaborados pelo autor que servem de ilustração, comprovação e ilustração.

\end{document}

\end{document}